\section{Thesis Structure}
\textbf{Chapter 2} introduces the background information related to the following areas: standard radar algorithms, monopulse angle estimation, multi-target tracking, and parallel processing.

\textbf{Chapter 3} conducts and extensive literature review of both foundational and state of the art approaches to performant radar systems under high SWaP constraints. The methodologies of works with similar motivations to ours are critically evaluated for their effectiveness in meeting their aims. Key findings and gaps across works are summarised and an explanation for how this thesis fits in with the existing literature is provided. 

\textbf{Chapter 4} Explains the methodology by which we compared different radar processing pipeline designs for their suitability for use in detect and avoid systems, and compare different algorithms for their ability to detect small, faraway targets. This methodology will describe the algorithms, the implementations of these algorithms, the processes by which the implementations are profiled, and the processes by which the algorithms are compared for their impact on the overall sensitivity of the system. 

\textbf{Chapter 5} Presents the results associated with each of the experiments outlined in the methodology. The results are illustrated as tables and figures, with descriptions of the findings of each of the methodology's stages.

\textbf{Chapter 6} The key findings in the results are discussed. The results of different experiment stages are evaluated and explained. Recommendations are provided. The results are tied in with the problem and contributions.

\textbf{Chapter 7} The key findings are summarised and the experiment is examined for any areas for improvement or future work. The value of the work is assessed for its contributions to the industry partner, and wider contributions to the field of UAV DAA.

\todo{Do I include appendix content like WHS report?}